\appendix

\section[Toda relations from finite difference operators (after Anderson--Chen--Tseng)]{Toda relations from finite difference operators\\ (after Anderson--Chen--Tseng)}\label{sec:toda}

The proof of the Toda relations in~\cite{maeno.naito.sagaki:QKideal}
relies on Kato's earlier results~\cite{kato:loop}. For the quantum K~ring
$\QK_T(\Fl(n))$, there is another proof of these relations, using an argument
combining the results of Iritani, Milanov and
Tonita~\cite{Iritani:2013qka} with results of Givental
and Lee~\cite{givental_lee}. More precisely, it is shown in~\cite{Iritani:2013qka} that
the symbols of finite
difference operators annihilating the K-theoretic~$J$ function of a variety $X$ give relations in the quantum K ring of $X$. Givental and Lee's results from loc.\ cit.\ imply that the
K-theoretic~$J$ function of the complete flag variety is an eigenfunction of the
(finite difference) Toda Hamiltonians. This observation was made in the unpublished note~\cite{act:2017} of Anderson--Chen--Tseng, but removed from the published version of their paper. For the sake of completeness, we give a brief account below, and in the process fill in some of the details to make the argument complete.

We start with recalling the definition
of the K-theoretic $J$-function of the complete flag variety $X=\Fl(n)$.
Denote by $P_i= \wedge^i \cS_i$; it is known that these line bundles algebra generate~$\K_T(\Fl(n))$
over~$\K_T(\pt)$; see, e.g.,~\cite[Proposition~3.1]{GMSXZZ:Nakayama}. Furthermore,
the curve classes associated to the Novikov variables
$Q_i$ are dual to the classes
$c_1\bigl(\wedge^i \cS_i^*\bigr)$.
For a fixed effective (multi)degree~${d \in H_2(X)}$, let $L$ be the cotangent line bundle at the unique marked point on the moduli space $\overline{M}_{0,1}(X,d)$.
Let also $\phi^\alpha$, $\phi_\alpha$ denote Poincar\'e-dual bases for $\K_T(X)$. (For example,
one may take Schubert classes $\cO^w$, and their duals -- the ideal sheaves of the boundary of the opposite Schubert varieties.)
The small $J$-function of $X$, denoted by $J_X$, is defined by
\[
J_X(q)\coloneq (1-q)\prod_i P_i^{\frac{\ln(Q_i)}{\ln(q)}}\sum_{d,\alpha} Q^d\biggl\langle \frac{\phi_\alpha}{1-qL}\biggr\rangle_{0,1,d}\phi^\alpha .
\]
We will explain later the meaning and the effect of the factor
\smash{$P_i^{\ln(Q_i)/\ln(q)}$}. We also note that the presence
of the prefactors $(1-q)$
and \smash{$\prod_i P_i^{\ln(Q_i)/\ln(q)}$} varies in the literature.
Our description agrees with the one used by Givental and Lee in
\cite{givental_lee}, and corresponds to the function denoted by $\tilde{J}$ in~\cite{Iritani:2013qka}.

We recall some basics on the formalism of difference operators. Consider commuting
variables $q, x_1, \dots, x_n$, and define the difference operators
\[
T_i:= q^{x_i \partial_{x_i}} = \sum_{ k \ge 0}^\infty \frac{1}{k!}((\ln q)x_i \partial_{x_i} )^k.
\]
\big(More generally, for a differential operator $\mathfrak{f}$, one defines the $q$-difference operator $q^{\mathfrak{f}}={\rm e}^{(\ln q) \mathfrak{f}}=\sum_{j=0}^\infty \frac{1}{j!}((\ln q)\mathfrak{f})^j$.\big) Note that
\[
T_i \bigl(x_j^{\pm 1}\bigr) = \sum_{k=0}^\infty \frac{1}{k!} ( \ln(q) x_i \partial_{x_i} ) \bigl(x_j^{\pm 1}\bigr) = q^{\pm \delta_{ij}} x_j^{\pm 1} ,
\]
which explains the `difference operator' terminology. More generally, for any Laurent polynomial in commuting variables $x_i$, we have
\[
T_i f(x_1,\dots,x_i,\dots,x_n) = f(x_1,\dots,qx_i,\dots, x_n) ,
\]
i.e., $T_i$ is an automorphism of the Laurent polynomial ring
$\Z\bigl[q^{\pm 1}; x_1^{\pm 1}, \dots, x_{n}^{\pm 1}\bigr]$.
We use this expression to extend the definition of $T_i$ to any function in the indeterminates $q, x_1, \dots, x_n$.

Now consider the subring of Laurent polynomials
\[
\Z\bigl[q^{\pm 1}; Q_1^{\pm 1}, \dots, Q_{n-1}^{\pm 1}\bigr] \hookrightarrow \Z\bigl[q^{\pm 1}; x_1^{\pm 1}, \dots, x_{n}^{\pm 1}\bigr]
\]
obtained by sending \smash{$Q_i \mapsto q^{-1} \frac{x_{i+1}}{x_i}$}. The restriction of $T_i$ to this subring is given by
\begin{equation}\label{E: TiQ}
T_i = q^{-Q_{i}\partial_{Q_{i}}}q^{Q_{i-1}\partial_{Q_{i-1}}} ,
\end{equation}
where \smash{$q^{Q_i \partial_{Q_i}}$} are the difference operators on the subring in $Q_i$'s.

With that in mind, we can now explain the meaning of the factor \smash{$P^{\frac{\ln(Q_i)}{\ln(q)}}$}.
The difference operators \smash{$q^{Q_i \partial_{Q_i}}$} act on functions in $Q_i$'s, and
one calculates that
\[ q^{Q_i \partial_{Q_i}}\bigl(P^{\frac{\ln(Q_j)}{\ln(q)}}\bigr) =
P^{\frac{\ln(q^{\delta_{ij}} Q_j)}{\ln(q)}} = P^{\delta{ij}}
P^{\frac{\ln(Q_j)}{\ln(q)}} . \]
In other words, the factor \smash{$P^{\frac{\ln(Q_i)}{\ln(q)}}$} should be regarded as
a formal variable
which transforms according to the rule above under the difference operators.

The relations in the quantum K ring are given by the Hamiltonians of the finite difference
(or relativistic) Toda lattice. There is some ambiguity in the exact expressions for the
Toda Hamiltonians, since their construction depends on choices;
see, e.g.,~\cite[Remark~5]{givental_lee}. We follow here the approach from~\cite{gloI},
but we will also need to make some changes of variables,
in order to fit with the conventions in our main reference~\cite{givental_lee}.
For the convenience of the reader, we briefly included some of the details below.

The Hamiltonians of the $q$-deformed type $A$ Toda chain have the form
\begin{align}\label{eqn:hamil}
 H_k=\sum_{0= i_0<\dots<i_k\leq n} \prod_{l=1}^{k}\biggl(1-\frac{x_{i_l}}{x_{i_l-1}}\biggr)^{1-\delta_{i_l-i_{l-1},1}} \prod_{l=1}^{k} T_{i_l},\qquad k=1,\dots,n,
\end{align}
where $q$ and $x_i$ are commuting variables, and $T_i=q^{x_i\partial_{x_i}}$ is the $q$-difference operator above. It was proved in~\cite{gloI} that the operators $H_k$ are limits of Macdonald operators, and the latter are known to commute.
This implies that $H_k$ also commute.

As above, let $Q_i=q^{-1}{x_{i+1}}{x_i}^{-1}$, with $Q_0=Q_n=0$. Then, using
\eqref{E: TiQ}, one can rewrite~\eqref{eqn:hamil} as
\begin{align*}
 H_k=\sum_{0= i_0<\dots<i_k\leq n} \prod_{l=1}^{k}(1-qQ_{i_l-1})^{1-\delta_{i_l-i_{l-1},1}} \prod_{l=1}^{k} q^{-Q_{i_l}\partial_{Q_{i_l}}}q^{Q_{i_l-1}\partial_{Q_{i_l-1}}},\qquad k=1,\dots,n.
\end{align*}
Replacing $q$ by $q^{-1}$, we obtain
\begin{align*}%\label{E:Toda-ops-changed}
 \widehat{H}_k&{}= \sum_{0= i_0<\dots<i_k\leq n}\prod_{l=1}^{k}\bigl(1-q^{-1}Q_{{i_l}-1}\bigr)^{1-\delta_{i_l-i_{l-1},1}} \prod_{l=1}^{k} q^{Q_{i_l}\partial_{Q_{i_l}}-Q_{i_l-1}\partial_{Q_{{i_l}-1}}}\\
 &{}=\sum_{0= i_0<\dots<i_k\leq n} \prod_{l=1}^{k} q^{Q_{i_l}\partial_{Q_{i_l}}-Q_{i_l-1}\partial_{Q_{{i_l}-1}}}\prod_{l=1}^{k}(1-Q_{{i_l}-1})^{1-\delta_{i_l-i_{l-1},1}}
 ,\qquad k=1,\dots,n .
 \end{align*}
\begin{Remark} The substitutions above
ensure that the first Hamiltonian $\widehat{H}_1$ agrees
with the one used in~\cite{givental_lee}. The substitution chosen in~\cite{act:2017}
produces similar
operators, but with the $q$-shifts and the Novikov terms in the opposite order.
\end{Remark}
The following key result of Givental and Lee~\cite[Theorem~2]{givental_lee} shows that the $J$ function is an eigenfunction for $J_{\Fl(n)}$.

\begin{Theorem}[Givental--Lee]\label{thm:GLeemain}
 $\widehat{H}_1J_{\Fl(n)}=\mathbb{C}^n J_{\Fl(n)}$.
\end{Theorem}

We also need the following lemma of Givental--Lee~\cite{givental_lee}.

\begin{Lemma}[{\cite[p.~9]{givental_lee}}]\label{lemma:modQ}
Let $D$ be a difference operator commuting with \smash{$\widehat{H}_1$}. Then, if $J$ is an eigenfunction of $D$ modulo $Q$, then $J$ is an eigenfunction of $D$ whose eigenvalue is the same as the one modulo $Q$.
\end{Lemma}

From this, we deduce that $J_{\Fl(n)}$ is an eigenfunction of the higher Toda Hamiltonians, using their commutativity with \smash{$\widehat{H}_1$} and by computing their eigenvalues modulo $Q$.

\begin{Corollary}\label{cor:Hk-eigenval} For any $1 \le k \le n$, the following holds:
\[ \widehat{H}_kJ_{\Fl(n)}=\wedge^k(\mathbb{C}^n) J_{\Fl(n)} . \]
\end{Corollary}
\begin{proof} The case $k=1$ is Theorem~\ref{thm:GLeemain}. Suppose $2 \le k \le n$. Since \smash{$\widehat{H}_k$} commutes with \smash{$\widehat{H}_1$}, we need only verify that $J_{\Fl(n)}$ is an eigenfunction of \smash{$\widehat{H}_k$} modulo $Q$, thanks to Lemma~\ref{lemma:modQ}.

To this end, we first observe
\begin{align*}
\widehat{H}_k J_{\Fl(n)}&{}=\widehat{H}_k\biggl((1-q)\prod_i P_i^{\frac{ln(Q_i)}{ln(q)}}\biggr) + o(Q_i) \\
&{}=
\sum_{0= i_0<\dots<i_k\leq n}\prod_{l=1}^k \frac{P_{i_l}}{P_{i_l-1}}\biggl((1-q)\prod_i P_i^{\frac{ln(Q_i)}{ln(q)}}\biggr)+o(Q_i) .
\end{align*}
Thus, modulo $Q$, we have the eigenvalue equation
\begin{align*}
\widehat{H}_k J_{\Fl(n)}&{}=\sum_{0= i_0<\dots<i_k\leq n}
\prod_{l=1}^k \frac{P_{i_l}}{P_{i_l-1}}J_{\Fl(n)}\\
&{}=
e_k\biggl(\frac{P_1}{P_0},\frac{P_2}{P_1},\dots,\frac{P_{n}}{P_{n-1}}\biggr)J_{\Fl(n)}=\wedge^k(\mathbb{C}^n)J_{\Fl(n)} .
\tag*{\qed}
\end{align*}
\renewcommand{\qed}{}
\end{proof}

We now use~\cite[Proposition~2.12]{Iritani:2013qka} which shows that the symbols of
Toda Hamiltonians give relations in quantum K theory.

\begin{Theorem}[Iritani--Milanov--Tonita]\label{thm:IMT}
Let $D=D\bigl(q^{Q_i\partial_{Q_i}},q,Q,\Lambda_i\bigr)$ be any $q$-difference operator
with coefficients in $\K_T(\pt)[q^{\pm 1}][\![Q_i]\!]$, such that it is regular
at $q=1$. Then
\[DJ_X=0\implies D\bigl(\widehat{P}_i,1,Q,\Lambda_i\bigr)=0 \in \QK_T(X) .\]
\end{Theorem}

\begin{Remark}
The result of Iritani, Milanov and Tonita is stated non-equivariantly, and for the
{\em big} quantum K ring and the corresponding big $J$ function. However, an inspection of their proof shows that it works in the equivariant situation as well.
Furthermore, if one starts with the small quantum K ring, then all arguments extend
to that situation, and the result also holds for the small quantum K ring and the small $J$ function. For further details, see~\cite{huq2024relations}.
\end{Remark}
One subtle point is that $\widehat{P}_i$ is a certain $Q$-deformation
of the line bundle $P_i$: it is the restriction to the small
quantum K ring
of an operator denoted by
$A_{i,{\rm com}}$ in~\cite[Corollary~2.9]{Iritani:2013qka}, which
arises as a solution to a certain Lax-type equation.
However, results of both Anderson, Chen, Tseng, and Iritani
in~\cite[Lemma~6]{anderson2022finite}, and also
by Kato in~\cite[Theorem~1.35]{kato:loop}
show that in fact no quantization is needed.

\begin{Proposition}\label{thm:ACTI-Kato} For the flag variety $\Fl(n)$,
$\widehat{\det(\mathcal{S}_i)}=\det(\mathcal{S}_i)$.
\end{Proposition}
We note in passing that an analogue of Proposition~\ref{thm:ACTI-Kato} holds for any homogeneous
space~$G/P$, but we do not need this generality here.

Combining Theorem~\ref{thm:IMT} and Proposition~\ref{thm:ACTI-Kato} with Corollary~\ref{cor:Hk-eigenval}
yields the following corollary.
\begin{Corollary}%\label{cor:QKToda}
The following identities hold
in $\QK_T(\Fl(n))$:
\begin{equation*}
 \sum_{0= i_0<\dots<i_k\leq n}\ \prod_{l=1}^{k}\frac{P_{i_l}}{P_{i_{l}-1}}(1-Q_{i_l-1})^{1-\delta_{i_l-i_{l-1},1}}=\wedge^k \C^n,\qquad k=1,\dots,n,
\end{equation*}
where $P_0=P_n=1$.
\end{Corollary}

\begin{Remark}
Theorem 4.9 of~\cite{koroteev} gives a presentation of the quasimap quantum $K$-ring of $T^*Fl$ whose limit to the is described in Theorem 5.5. The relations are based on the trigonometric Ruijsenaars--Schneider model. After further taking into account a restriction from $\GL_n$ to $\SL_n$, the `Toda limit' recovers the relations in this paper. We are grateful to P. Koroteev who explained this procedure to us.
\end{Remark}

